\section{Posicionamiento y estrategia de SSI}

\subsection{Straight-Shot Super Intelligence?}

El plan original de SSI era "ir directo a la superinteligencia" (straight-shot superintelligence), sin pasar por la competencia de mercado. Ilya expuso consideraciones tanto a favor como en contra:

\textbf{A favor de ir directo}:
\begin{itemize}
    \item No verse afectado por las presiones cotidianas de la competencia de mercado, pudiendo concentrarse en la investigación
    \item Evitar verse obligado a tomar trade-offs difíciles
\end{itemize}

\textbf{En contra de ir directo}:
\begin{itemize}
    \item Dejar que el mundo vea que una IA poderosa tiene valor: "communicate the AI, not the idea"
    \item Si el horizonte de tiempo es largo, un modo puramente de investigación no es sostenible
    \item La mejor IA se usa y mejora a través del uso (por analogía con la mejora de la seguridad de los aviones, que proviene de los vuelos reales)
\end{itemize}

\subsection{La escasez de ideas}

\begin{importantbox}{If ideas are so cheap, how come no one's having any ideas?}
Uno de los legados de la era del scaling es que hay más empresas que ideas. Todos están haciendo lo mismo (RL on pre-trained models), y las ideas realmente diferenciadoras son escasas. SSI afirma tener una vía técnica distinta (centrada en la generalización y la comprensión), que tendrá un valor considerable si estas ideas resultan correctas.
\end{importantbox}

\subsection{La salida del cofundador}

Ilya respondió brevemente al episodio de la salida de un cofundador para incorporarse a Meta: en ese momento SSI estaba recaudando fondos con una valoración de 32,000 millones de dólares, y Meta hizo una oferta de adquisición. Ilya la rechazó, pero el excofundador optó por incorporarse a Meta (y fue también la única persona que dejó SSI para irse a Meta).

\subsection{Resumen de la sección}

SSI se posiciona como "la empresa de la era de la investigación", con una vía técnica diferenciadora centrada en la generalización. La estrategia original de "ir directo a la superinteligencia" podría ajustarse por consideraciones de horizonte de tiempo y de valor del despliegue.

